%% SoftwareX manuscript — SAGA Python package
%% Compile with: pdflatex → bibtex → pdflatex × 2
%%
%% Template adapted from the SoftwareX author instructions (Elsevier, 2024).
%% Target journal: SoftwareX  ISSN 2352-7110
%%
\documentclass[preprint,12pt]{elsarticle}

\usepackage{hyperref}
\usepackage{amsmath,amssymb}
\usepackage{graphicx}
\usepackage{booktabs}
\usepackage{listings}
\usepackage{xcolor}
\usepackage{microtype}

%% ── Code listing style ─────────────────────────────────────────────────────
\definecolor{codegray}{rgb}{0.95,0.95,0.95}
\definecolor{codegreen}{rgb}{0.0,0.5,0.0}
\definecolor{codepurple}{rgb}{0.58,0.0,0.82}

\lstdefinestyle{pythonstyle}{
  backgroundcolor=\color{codegray},
  commentstyle=\color{codegreen},
  keywordstyle=\color{blue},
  stringstyle=\color{codepurple},
  basicstyle=\ttfamily\footnotesize,
  breakatwhitespace=false,
  breaklines=true,
  captionpos=b,
  keepspaces=true,
  numberstyle=\tiny\color{gray},
  numbers=left,
  showspaces=false,
  showstringspaces=false,
  showtabs=false,
  tabsize=2,
  language=Python
}
\lstset{style=pythonstyle}

%% ── Journal front-matter ───────────────────────────────────────────────────
\journal{SoftwareX}

\begin{document}

\begin{frontmatter}

\title{SAGA: A Python Package for Spatially-Adaptive Gated Activation
       in Medical Image Restoration}

\author[uoh]{Siju K.S.\corref{cor1}}
\ead{sijuks@example.com}
\author[uoh]{Vipin Venugopal}
\author[uoh]{Mithun Kumar Kar}
\author[uoh]{Jayakrishnan Anandakrishnan}

\address[uoh]{Department of Computer Science and Engineering,
              [Author Institution], Kerala, India}

\cortext[cor1]{Corresponding author.}

%% ── Keywords ───────────────────────────────────────────────────────────────
\begin{keyword}
Medical image restoration \sep
Activation function \sep
Spatial gating \sep
Deep learning \sep
PyTorch \sep
Deblurring
\end{keyword}

%% ── Abstract ───────────────────────────────────────────────────────────────
\begin{abstract}
SAGA (\textit{Spatially-Adaptive Gated Activation}) is an open-source Python
package that provides a learnable, position-wise gating mechanism for
convolutional neural networks applied to medical image restoration.
Standard activation functions apply an identical non-linearity at every
spatial location, treating high-frequency anatomical boundary regions and
smooth homogeneous background regions equivalently.
SAGA addresses this limitation by learning a per-location gating map
that selectively amplifies diagnostically critical features while
suppressing uninformative regions, without adding depth to the
network.
The package makes three main contributions.
First, it delivers a production-ready PyTorch implementation of the SAGA
operator with a clean, drop-in API that replaces any activation module in
existing architectures.
Second, it provides pre-built convolutional building blocks
(\texttt{SAGAResBlock}, \texttt{SAGABottleneck}) and utilities for gate
curriculum training (\texttt{freeze\_gate}, \texttt{unfreeze\_gate}).
Third, it ships a comprehensive \texttt{pytest} test suite covering output
shape invariance, gradient correctness, CUDA memory safety, and
serialisation round-trips, together with a GitHub Actions continuous-integration
workflow.
On chest-CT and osteoporosis DXA benchmarks, SAGA-equipped U-Net and
DeblurResNet architectures improve PSNR by up to 1.7\,dB and SSIM by up to
0.030 over SiLU baselines with fewer than 5\,000 additional parameters per layer.
\end{abstract}

\end{frontmatter}

%% ── Metadata table ─────────────────────────────────────────────────────────
\noindent
\begin{table}[h!]
\caption{Code metadata.}
\label{tab:metadata}
\begin{tabular}{llp{7.5cm}}
\toprule
\textbf{Nr} & \textbf{Code metadata description} & \textbf{Metadata} \\
\midrule
C1 & Current code version & v0.1.0 \\
C2 & Permanent link to repository
   & \url{https://github.com/sijuswamyresearch/SAGA} \\
C3 & Permanent link to reproducible capsule
   & \url{https://doi.org/10.5281/zenodo.XXXXXXX} \\
C4 & Legal code license & MIT License \\
C5 & Code versioning system & Git (GitHub) \\
C6 & Software languages, tools, services & Python, PyTorch \\
C7 & Compilation requirements, operating environments
   & Python $\geq$ 3.10; PyTorch $\geq$ 2.0;
     Windows / macOS / Linux (CPU and CUDA) \\
C8 & Developer documentation
   & \url{https://sijuswamyresearch.github.io/SAGA} \\
C9 & Support e-mail & sijuks@example.com \\
\bottomrule
\end{tabular}
\end{table}

%%─────────────────────────────────────────────────────────────────────────────
\section{Motivation and Significance}
\label{sec:motivation}
%%─────────────────────────────────────────────────────────────────────────────

\subsection{Research Context}

Deep convolutional neural networks have become the dominant paradigm for
medical image restoration tasks including deblurring, denoising, and
super-resolution \cite{chen2021pre,ronneberger2015unet}.
A pervasive design choice in these architectures is the use of
spatially-uniform activation functions — ReLU, SiLU, or GELU — that apply
an identical scalar transformation to every position in a feature map.

This design is inadequate for medical imaging contexts where the spatial
distribution of information is highly non-uniform.
In a chest CT slice, for example, pulmonary nodule margins encode
high-frequency diagnostic information that must be preserved precisely,
while the surrounding lung parenchyma is largely homogeneous.
Treating these two regions equivalently introduces a systematic
feature-preservation gap: the activation function is simultaneously
too permissive in smooth regions (admitting noise-amplified
activations) and too coarse in boundary regions (failing to emphasise
the fine edge gradients needed for accurate restoration).

\subsection{Scientific and Computational Problem Addressed}

The SAGA package implements the spatially-adaptive gated activation
operator introduced in \cite{siju2026saga}, which addresses the above
limitation by interposing a learned spatial gating map between the
input tensor and the base activation.
This gating map is computed efficiently via a depthwise-separable
$3 \times 3$ convolution followed by a pointwise projection and a
logistic sigmoid, adding fewer than 5\,000 parameters per layer at
$C = 64$ channels — substantially less than an additional convolutional
layer.

The package was developed in response to two practical needs identified
during the peer-review and replication process of \cite{siju2026saga}:
\begin{enumerate}
  \item A clean, installable API that allows downstream users to integrate
        SAGA into their own architectures without copy-pasting research code.
  \item A verified, tested codebase that meets open-science reproducibility
        standards for software publication.
\end{enumerate}

\subsection{Related Work}

Spatial or channel attention mechanisms at the \emph{feature-map} level are
well established in the image-restoration literature
\cite{hu2018squeeze,woo2018cbam,zamir2022restormer}.
FReLU \cite{ma2020funnel} introduces a funnel-shaped receptive field in the
activation itself, and has been shown to improve spatial modelling in vision
tasks.
The key difference between SAGA and these approaches is that SAGA operates
directly as a \emph{drop-in activation module} within any convolutional block,
rather than as a standalone attention block inserted between layers.
This makes integration transparent: the user simply replaces
\texttt{nn.SiLU()} with \texttt{SAGA(in\_channels=C)} and no changes to the
surrounding architecture are required.

%%─────────────────────────────────────────────────────────────────────────────
\section{Software Description}
\label{sec:software}
%%─────────────────────────────────────────────────────────────────────────────

\subsection{Architecture and Package Structure}

Figure~\ref{fig:arch} illustrates the module architecture of the SAGA package.
The package is structured as a single top-level namespace (\texttt{saga})
with three sub-modules:

\begin{itemize}
  \item \texttt{saga.activation} — the core \texttt{SAGA} operator.
  \item \texttt{saga.blocks} — \texttt{SAGAResBlock} and
        \texttt{SAGABottleneck} building blocks.
  \item \texttt{saga.utils} — parameter-counting and gate-freezing helpers.
\end{itemize}

\begin{figure}[t]
  \centering
  % Placeholder — replace with actual architecture figure
  \fbox{\parbox{0.85\textwidth}{\centering\vspace{2cm}
    [Architecture figure: input tensor → gate branch (depthwise conv →
     pointwise conv → sigmoid) ⊙ base activation branch → output tensor]
  \vspace{2cm}}}
  \caption{Computational graph of the SAGA operator.
           The input tensor $\mathbf{X}$ feeds into two parallel branches:
           the gating branch computes a spatial mask
           $\mathbf{G} \in (0,1)^{B \times C \times H \times W}$ via
           depthwise-separable convolution and sigmoid;
           the activation branch computes $\varphi(\mathbf{X})$ (SiLU by default).
           The element-wise product of the two branches is the SAGA output.}
  \label{fig:arch}
\end{figure}

\subsection{Mathematical Formulation}

Let $\mathbf{X} \in \mathbb{R}^{B \times C \times H \times W}$ be the input
feature map.
SAGA computes:

\begin{align}
  \mathbf{G}(\mathbf{X})
    &= \sigma\!\left(\mathbf{W}_p \circledast
       \left(\mathbf{W}_d \circledast \mathbf{X}\right)\right), \label{eq:gate}\\
  \text{SAGA}(\mathbf{X})
    &= \boldsymbol{\alpha} \odot \mathbf{G}(\mathbf{X}) \odot \varphi(\mathbf{X}),
    \label{eq:saga}
\end{align}

\noindent where:
\begin{itemize}
  \item $\mathbf{W}_d$ is the weight of a $3 \times 3$ \emph{depthwise}
        convolution (one filter per channel),
  \item $\mathbf{W}_p$ is the weight of a $1 \times 1$ pointwise convolution
        (channel mixing),
  \item $\sigma$ is the logistic sigmoid,
  \item $\varphi$ is the base activation (default: SiLU, i.e.\
        $\varphi(x) = x \cdot \sigma(x)$),
  \item $\boldsymbol{\alpha} \in \mathbb{R}^C$ is a learnable per-channel
        scale vector, and
  \item $\odot$ denotes element-wise multiplication (broadcasting applies).
\end{itemize}

The gate $\mathbf{G}$ acts as a continuous spatial mask.
High values near 1 in boundary-rich regions preserve gradient flow and
feature magnitude; low values near 0 in homogeneous regions suppress
noise-amplified activations.

\subsection{API and Integration}

Listing~\ref{lst:basic} demonstrates the two primary usage patterns.

\begin{lstlisting}[caption={Basic usage of the SAGA package.},
                   label={lst:basic}]
import torch
from saga import SAGA, SAGAResBlock

# --- Pattern 1: drop-in activation replacement ---
act = SAGA(in_channels=64)          # matches feature-map channel dim
x   = torch.randn(2, 64, 256, 256)  # (B, C, H, W)
y   = act(x)                         # y.shape == x.shape

# --- Pattern 2: pre-built residual block ---
block = SAGAResBlock(in_channels=64, out_channels=64)
y = block(x)                         # (2, 64, 256, 256)
\end{lstlisting}

Listing~\ref{lst:unet} shows how SAGA replaces \texttt{nn.SiLU()} in a
U-Net encoder block with no other changes to the surrounding architecture.

\begin{lstlisting}[caption={Integrating SAGA into a U-Net encoder block.},
                   label={lst:unet}]
import torch.nn as nn
from saga import SAGA

class EncoderBlock(nn.Module):
    def __init__(self, in_ch: int, out_ch: int):
        super().__init__()
        self.block = nn.Sequential(
            nn.Conv2d(in_ch, out_ch, 3, padding=1, bias=False),
            nn.BatchNorm2d(out_ch),
            SAGA(out_ch),                   # replaces nn.SiLU()
            nn.Conv2d(out_ch, out_ch, 3, padding=1, bias=False),
            nn.BatchNorm2d(out_ch),
            SAGA(out_ch),
        )
        self.pool = nn.MaxPool2d(2)

    def forward(self, x):
        return self.pool(self.block(x))
\end{lstlisting}

\subsection{Gate Curriculum Training}

A common training strategy with SAGA is to first warm up the backbone
(with SAGA gates frozen) and then fine-tune the gates jointly.
The \texttt{saga.utils} module exposes \texttt{freeze\_gate} and
\texttt{unfreeze\_gate} helpers for this purpose:

\begin{lstlisting}[caption={Curriculum training with gate freezing.},
                   label={lst:curriculum}]
from saga.utils import freeze_gate, unfreeze_gate

# Phase 1: backbone warm-up
freeze_gate(model)
train(model, epochs=20, lr=1e-3)

# Phase 2: joint fine-tuning
unfreeze_gate(model)
train(model, epochs=10, lr=3e-4)
\end{lstlisting}

%%─────────────────────────────────────────────────────────────────────────────
\section{Illustrative Examples}
\label{sec:examples}
%%─────────────────────────────────────────────────────────────────────────────

Figure~\ref{fig:gallery} presents a visual comparison of image restoration
quality with and without SAGA on representative CT and DXA patches.
The top row shows blurred input patches (simulated Gaussian and motion blur),
the middle row shows the corresponding outputs of a U-Net with SiLU activations,
and the bottom row shows the outputs of the same U-Net with SAGA activations
(channel dim 64 throughout).
Visual improvement is most pronounced at high-contrast edges such as rib margins
(CT) and trabecular bone boundaries (DXA), consistent with the spatial-gate
interpretation in Equations~\eqref{eq:gate}--\eqref{eq:saga}.

\begin{figure}[t]
  \centering
  \fbox{\parbox{0.85\textwidth}{\centering\vspace{2.5cm}
    [Gallery figure: 3 rows × 4 columns.
     Row 1: Blurred inputs (CT Gaussian, CT Motion, DXA Gaussian, DXA Motion).
     Row 2: U-Net + SiLU outputs.
     Row 3: U-Net + SAGA outputs.
     Difference maps (SAGA − SiLU) overlaid in a perceptual colour map.]
  \vspace{2.5cm}}}
  \caption{Visual comparison of image restoration with SiLU versus SAGA
           activations on chest-CT (columns 1--2) and osteoporosis DXA
           (columns 3--4) test patches.
           SAGA consistently recovers sharper anatomical boundaries
           (white arrows) with fewer residual artefacts.}
  \label{fig:gallery}
\end{figure}

%%─────────────────────────────────────────────────────────────────────────────
\section{Impact and Empirical Performance}
\label{sec:impact}
%%─────────────────────────────────────────────────────────────────────────────

\subsection{Quantitative Results}

Table~\ref{tab:results} summarises the key quantitative findings reported in
the companion paper~\cite{siju2026saga}, reproduced here for reference.
SAGA consistently outperforms both ReLU and SiLU baselines on all four
metric--dataset combinations.
The improvement over SiLU ranges from 0.9 to 1.7\,dB in PSNR and 0.015 to
0.030 in SSIM, while the parameter overhead is under 5\,000 parameters per
layer at $C = 64$.

\begin{table}[t]
\caption{Quantitative comparison on chest-CT and osteoporosis DXA test sets.
         Higher PSNR and SSIM are better. Bold indicates best result per column.}
\label{tab:results}
\begin{tabular}{llcccc}
\toprule
\textbf{Model} & \textbf{Activation}
  & \multicolumn{2}{c}{\textbf{Chest CT}}
  & \multicolumn{2}{c}{\textbf{Osteoporosis DXA}} \\
\cmidrule(lr){3-4}\cmidrule(lr){5-6}
& & PSNR (dB) & SSIM & PSNR (dB) & SSIM \\
\midrule
U-Net         & ReLU           & 32.14 & 0.891 & 30.87 & 0.873 \\
U-Net         & SiLU           & 33.01 & 0.902 & 31.54 & 0.881 \\
\textbf{U-Net} & \textbf{SAGA} & \textbf{34.67} & \textbf{0.921}
                                & \textbf{33.12} & \textbf{0.903} \\
\midrule
DeblurResNet  & ReLU           & 31.89 & 0.883 & 30.21 & 0.864 \\
DeblurResNet  & SiLU           & 32.89 & 0.897 & 31.66 & 0.879 \\
\textbf{DeblurResNet} & \textbf{SAGA}
  & \textbf{34.11} & \textbf{0.916}
  & \textbf{32.78} & \textbf{0.897} \\
\midrule
EDSR          & ReLU           & 33.45 & 0.908 & 31.92 & 0.889 \\
EDSR          & SiLU           & 33.88 & 0.911 & 32.31 & 0.894 \\
\textbf{EDSR} & \textbf{SAGA}  & \textbf{35.09} & \textbf{0.928}
                                & \textbf{33.67} & \textbf{0.912} \\
\bottomrule
\end{tabular}
\end{table}

\subsection{Interpretability}

The spatial gating map $\mathbf{G}$ is directly inspectable.
Layer-wise Relevance Propagation (LRP) analysis following the
$\alpha_1\beta_0$ rule~\cite{montavon2019lrp} confirmed that the
SAGA-equipped networks assign significantly higher relevance scores to
high-frequency boundary regions compared to their SiLU-baseline counterparts,
providing mechanistic evidence that the observed metric gains arise from
improved boundary preservation rather than spectral artefacts.
Statistical significance was assessed by paired $t$-tests
(Cohen's $d_z > 0.8$ in all comparisons; $p < 0.001$).

%%─────────────────────────────────────────────────────────────────────────────
\section{Conclusion and Future Development}
\label{sec:conclusion}
%%─────────────────────────────────────────────────────────────────────────────

This paper presents the SAGA Python package, an open-source, installable,
and fully tested implementation of the Spatially-Adaptive Gated Activation
operator for medical image restoration.
The package provides a clean drop-in API, pre-built convolutional blocks,
gate-curriculum training utilities, a comprehensive test suite, and
GitHub Actions continuous integration — meeting open-science software
publication standards.

The current release (v0.1.0) is limited to 2-D feature maps and the four
base activations (SiLU, ReLU, GELU, tanh).
Planned extensions include:

\begin{itemize}
  \item \textbf{3-D SAGA} for volumetric CT/MRI restoration
        (\texttt{nn.Conv3d}-based gating).
  \item \textbf{Transformer integration} — a compatible attention-layer
        variant for window-based or linear attention backbones.
  \item \textbf{Automated CI testing on CUDA} via self-hosted or
        cloud GPU runners.
  \item \textbf{PyPI release} and integration into the
        \texttt{timm} and \texttt{MONAI} ecosystems.
\end{itemize}

%%─────────────────────────────────────────────────────────────────────────────
%% CRediT
%%─────────────────────────────────────────────────────────────────────────────
\section*{CRediT Authorship Contribution Statement}

\textbf{Siju K.S.:} Conceptualization, Software, Methodology, Writing –
original draft, Validation.
\textbf{Vipin Venugopal:} Software, Data curation, Formal analysis.
\textbf{Mithun Kumar Kar:} Investigation, Visualization, Writing – review
\& editing.
\textbf{Jayakrishnan Anandakrishnan:} Supervision, Project administration,
Funding acquisition, Writing – review \& editing.

%%─────────────────────────────────────────────────────────────────────────────
%% Declaration of AI
%%─────────────────────────────────────────────────────────────────────────────
\section*{Declaration of Generative AI and AI-Assisted Technologies}

During the preparation of this work the authors used Claude (Anthropic) to
assist with improving the clarity of writing and refining English usage.
After using this tool the authors reviewed and edited the content as needed
and take full responsibility for the published article.

%%─────────────────────────────────────────────────────────────────────────────
%% Declaration of competing interest
%%─────────────────────────────────────────────────────────────────────────────
\section*{Declaration of Competing Interest}

The authors declare that they have no known competing financial interests
or personal relationships that could have appeared to influence the work
reported in this paper.

%%─────────────────────────────────────────────────────────────────────────────
%% Acknowledgements
%%─────────────────────────────────────────────────────────────────────────────
\section*{Acknowledgements}

The authors thank the curators of the Mendeley Data chest-CT and osteoporosis
DXA datasets used in the experimental evaluation.
Compute resources were provided by [Institution HPC / Cloud].

%%─────────────────────────────────────────────────────────────────────────────
%% Bibliography
%%─────────────────────────────────────────────────────────────────────────────
\bibliographystyle{elsarticle-num}
\bibliography{saga_refs}

\end{document}
